\section{Introduction}\label{sec:introduction}

A projected entangled-pair state (PEPS) represents a many-body vector by local
tensors joined along the edges of its interaction lattice. The central
approximation question is whether ground states of gapped local Hamiltonians
admit such representations with modest virtual dimensions. An entropy area
law controls every bipartite cut, but a global vector approximation requires
more: the descriptions at different cuts must be compatible, their errors
must remain small through the entire lattice, and the final tensor network
must use the prescribed edges.

We prove a uniform polynomial bound for square lattices with arbitrary
bounded on-site and nearest-neighbor interactions. The only spectral
assumption is a unique ground vector and a positive gap for the full
Hamiltonian. We use the entropy area law and the collar estimates of
\cite{RepoAreaLaw}, always with their full-system hypotheses. The new argument
establishes the stronger representation conclusion.

\subsection{Statement of the result}

For an integer $L\ge2$, let $\Lambda_L=\{1,\ldots,L\}^2$, with the
nearest-neighbor square-lattice edge set $E(\Lambda_L)$ and open boundary
conditions. Put
\[
 \mathcal H_{\Lambda_L}=\bigotimes_{v\in\Lambda_L}\mathbb C^q.
\]
A PEPS on this grid is the contraction obtained by placing
$\sum_{a=1}^{D_e}|a\rangle\otimes|a\rangle$ on each edge $e$ and applying
at each vertex an arbitrary linear map
\[
 A_v:\bigotimes_{e\ni v}\mathbb C^{D_e}\longrightarrow\mathbb C^q.
\]
Its maximum bond dimension is $\max_eD_e$. The maps need not be isometries,
and the resulting vector need not be normalized.

\begin{theorem}\label{thm:main}
Fix an integer $q\ge2$ and constants $J,\Delta>0$. There are constants
$C,c>0$, depending only on $q,J,\Delta$, with the following property.
For every $L\ge2$, consider any Hermitian Hamiltonian
\begin{equation}\label{eq:model}
 H=\sum_{v\in\Lambda_L}h_v+\sum_{e\in E(\Lambda_L)}h_e,
 \qquad \norm{h_v},\norm{h_e}\le J,
\end{equation}
where each summand is supported on its designated site or edge.
Suppose that $H$ has a unique normalized ground vector $\Omega$, up to
phase, with ground energy $E_0$, and
\begin{equation}\label{eq:global-gap}
 H-E_0I\ge\Delta\bigl(I-|\Omega\rangle\langle\Omega|\bigr).
\end{equation}
Then there is a nonzero PEPS $\Phi$ on the same grid with maximum bond
dimension at most $CL^c$ such that
\begin{equation}\label{eq:target-error}
 \min_{\theta\in\mathbb R}
 \left\|\frac{\Phi}{\norm{\Phi}}-e^{i\theta}\Omega\right\|
 \le L^{-1}.
\end{equation}
\end{theorem}

Theorem~\ref{thm:main} gives an affirmative resolution of the polynomial
PEPS approximation assertion for the stated open-square Hamiltonian class.
The relation to other formulations of the PEPS conjecture is explained below. No gap for restrictions to subregions,
frustration freeness, commutativity, translation invariance, or path from a
product state is assumed. The error in \eqref{eq:target-error} concerns the
entire vector. The theorem makes no claim that its tensors can be found or
contracted efficiently; their entries may depend arbitrarily on $H$ and
$\Omega$.

\subsection{History and significance}
Projected entangled-pair states extend matrix product states to
higher-dimensional lattices by contracting virtual pairs through local
maps~\cite{VerstraeteCirac2004PEPS}. Their local tensor structure makes them
natural candidates for describing ground states whose entanglement is
concentrated near region boundaries. In one dimension, quantitative
Schmidt-tail truncation estimates connect entanglement spectra with global
MPS approximation~\cite{VerstraeteCirac2006MPS}. Hastings established the
entropy area law for unique ground states of uniformly gapped spin chains
with bounded finite-range interactions and obtained polynomial-dimension
MPS approximations~\cite{Hastings2007}. The higher-dimensional question
requires a representation compatible across many intersecting cuts.

There are several distinct PEPS approximation questions. The global
formulation in \cite[Section~2.2, Questions~7--8]{CiracGarreRubioPerezGarcia2019}
asks whether a spectral gap ensures polynomial bond dimension at
inverse-polynomial Hilbert-space norm error. That formulation uses
translation-invariant finite-range Hamiltonians on a torus.
Theorem~\ref{thm:main} addresses the corresponding assertion for arbitrary
bounded on-site and nearest-neighbor interactions on an open square.
Local-observable formulations instead control expectations on fixed-size
regions; the stronger structural formulation in
\cite[Conjecture~2]{SchwarzBuerschaperEisert2017} also requires an injective
approximant with a gapped parent Hamiltonian. The theorem here concerns the
full vector and imposes no such structural requirements on the approximant.

A positive local approximation result is due to Huang
\cite[Theorem~3]{Huang2026Local}. For states satisfying a R\'enyi entropy
area law of fixed order $0<\alpha<1$ across every rectangle, he obtains
PEPS of bond dimension $e^{O(1/\delta)}$, independent of lattice size,
that approximate norm-one local-observable expectations to error $\delta$.
The observable support scale is held fixed in this bound. The R\'enyi
hypothesis is stronger than a von Neumann entropy area law, and the
conclusion controls local expectations rather than global vector error.

An entropy bound and a representation theorem are different statements.
For arbitrary chain-state families, bounded von Neumann block entropy
does not ensure efficient global MPS approximation
\cite{SchuchWolfVerstraeteCirac2008}. In higher dimensions, Ge and Eisert
constructed families satisfying strong area laws that have no efficient
classical tensor-network descriptions
\cite[Definition~7 and Corollary~9]{GeEisert2016}. Their description model
includes a bound on the classical complexity of tensor entries; it should
be distinguished from a bound on virtual dimension with arbitrary complex
entries. These examples concern general states, rather than the uniformly
gapped ground-state class in Theorem~\ref{thm:main}. They explain why one
must control compatible operations and their accumulated error, in addition
to the spectra of separate marginals.

Thermal tensor networks provide another approach. Hastings constructed
local spectral filters and tensor-network descriptions of Gibbs operators
\cite{Hastings2006Local}. Moln\'ar, Schuch, Verstraete, and Cirac refined the
thermal approximation and derived quasipolynomial-dimension ground-state
PEPS under a gap together with a low-energy density-of-states bound
\cite[Section~III.C]{MolnarSchuchVerstraeteCirac2015}. At inverse-polynomial
accuracy their bond bound is exponential in a squared logarithm of system
size. The low-energy counting condition is an additional hypothesis.
A different cooling and postselection construction approximates the unique
ground state of a gapped $d$-dimensional local Hamiltonian at the boundary
of a $(d+1)$-dimensional PEPS
\cite[Section~V]{SchuchWolfVerstraeteCirac2007Complexity}; transferring that
representation to the original lattice is a further problem.

Recent tensor-network results also distinguish the dimension and the
approximation task. Arad, Firanko, and Jain obtained two-dimensional
mutual-information bounds for locally gapped frustration-free systems and
polynomial-bond matrix product operator approximations to maximally mixed
ground states in one dimension
\cite[Corollaries~2.7 and~2.9]{AradFirankoJain2026}. Such entropy or
one-dimensional operator bounds do not provide the simultaneous
original-grid PEPS conclusion needed here.

Our analytic inputs are the entropy area law and the stronger collar
estimates of the companion \cite{RepoAreaLaw}. The present proof adapts its
nested-filter and two-family ideas, and then develops compatible changes
of ownership, compression of a distributed density network, and routing on
the prescribed grid. The representation theorem uses the full companion
interface recorded in Appendix~\ref{app:companion}, rather than the
cutwise entropy inequality alone. Its polynomial bond bound permits
arbitrary tensor entries and is an existence statement. The distinction
from efficient contraction is substantive: general PEPS contraction has
worst-case complexity obstructions
\cite{SchuchWolfVerstraeteCirac2007Complexity}, although that fact alone does
not classify the particular approximants constructed here.


\subsection{Proof strategy}

The main intermediate object is a distribution of the physical registers
among auxiliary parties. A party owns a collection of registers and may
perform arbitrary linear contractions on them privately. At the start,
one party holds the entire vector $\Omega$; at the end, each physical
site has its own party. Preparing the initial vector is permitted because
we seek a representation, rather than an algorithm for finding it.
The constraint is on communication: each party must participate in only
boundedly many operations with other parties over its whole lifetime.
Each such operation must also have a polynomial-size expansion built from
private contractions and normalized bipartite vectors or covectors. The
bipartite states themselves may have arbitrarily large Schmidt rank.

\paragraph{Distributing the state while preserving a reference vector.}
The parties are indexed by dyadic squares. We pass from the ownership
assignment on one scale to that on the next by giving each child square
its physical registers. Intermediate assignments need not be dyadic, and
we use additional independent copies of $\Omega$. Each copy has its own
registers and ownership assignment; its reference vector remains $\Omega$
in the original site coordinates, up to the small-patch encodings described
below. The operations approximate these reference vectors with an error
small enough to sum over the whole construction.

There are two analytic inputs. First, information between regions becomes
negligible when their separation tapers linearly toward finitely many
points but stays above a sufficiently large multiple of $\log L$
(Proposition~\ref{prop:angular}). Small mutual information lets the owner
of a separating buffer split its purification approximately into two
factors. This implements a change of owner for the interior region, or
an exchange of assignments between two copies, using private maps and
bipartite resources. To prove the information estimate, an ordered
partition into two tile families bounds the conditional-information cost
of adding a cell. Gaussian filtering with the full-system gap then resets
that cost to an inverse power of $L$ after each annular step.

Second, near a logarithmic-size square we can project with high accuracy
onto a sum of mutually orthogonal subspaces. Writing
$\mathcal H_A=\bigotimes_{v\in A}\mathbb C^q$ for the registers in
$A\subseteq\Lambda_L$, these subspaces have the form
\[
 \operatorname{span}\{v_{j\ell}\}_{\ell=1}^{d_j}
       \otimes\mathcal H_{\Lambda_L\setminus X_j},
 \qquad \sum_jd_j\le CL^c,
\]
where the $X_j$ are nested square samples within the specified patch and
$v_{j\ell}\in\mathcal H_{X_j}$ are orthonormal inside vectors
(Proposition~\ref{prop:patch}). The radius is allowed to depend on $j$.
Recording $(j,\ell)$ in a small register and replacing the corresponding
inside vector by a fixed product vector gives the encoding used near
junctions. The proof combines optimized nested filters, the gap, and the
entropy area law; it does not require an accurate truncation at a single
prescribed boundary. Section~\ref{sec:frames} uses these encodings to
express suitable changes on small patches in terms of pair states and
effects, even when their underlying Hilbert spaces are large.

The remaining geometric problem is to arrange that these operations
suffice. Along a long edge, our separation estimate may not justify
changing an old owner directly to a new one: a third owner can touch the
opposite side of the region being changed. Exchanging part of an auxiliary copy first supplies the
surrounding regions needed for successive changes of owner. Small patches
handle the finitely many endpoints and junctions, where the separation
would otherwise vanish. Section~\ref{sec:geometry} assembles these moves
so that a block party participates only in boundedly many nearby
repaintings at each of two adjacent scales. Its total participation in
operations with other parties is therefore bounded.

\paragraph{Compressing the density and returning to a vector.}
The distribution initially uses pair resources of unrestricted dimension.
Theorem~\ref{thm:compression} replaces them in the density calculation by
polynomial sums of product operators. To bound the error from
replacing a chosen set of resources, we expand only operations involving
the parties that hold their endpoints. Bounded lifetime participation
limits the resulting coefficient cost. A Gaussian
second-moment bound then controls the trace-norm error independently of
all private dimensions and resource Schmidt ranks. This produces an
operator tensor network with polynomial virtual dimensions on the party
graph. The compression theorem is independent of the planar construction.

The protocol retains norm close to one, so its density and the compressed
operator both approximate $|\Omega\rangle\langle\Omega|$ in absolute trace
norm. One product-basis column of the compressed operator gives a nonzero
ket approximation after normalization. Finally, the links join parties
at comparable nearby dyadic scales. Routing them along rows and columns
through their dyadic anchors has constant edge congestion, and therefore
preserves polynomial bond dimension on the original square grid.
Section~\ref{sec:assembly} proves these last steps and the main theorem.

Section~\ref{sec:preliminaries} fixes the conventions and states the area
law. Sections~\ref{sec:information} and~\ref{sec:patches} prove the two
analytic inputs above; Appendix~\ref{app:companion} records their exact
companion interfaces. All error powers, geometric constants, and
polynomial exponents are fixed before taking $L$ large, uniformly in the
individual Hamiltonian.
